\documentclass[english, parskip=full]{scrartcl}
\usepackage[utf8]{inputenc}  % Kodierung der Datei
\usepackage[english]{babel}  % Sprache auf Deutsch 
\usepackage{color}
\usepackage{graphicx, gensymb}
\usepackage{amsfonts, cancel, amssymb}
\usepackage{amsmath, amsthm, array, esint}
\usepackage{booktabs, multirow}  % schönere Tabellen
\usepackage{caption}  % Anpassung der Captions
\usepackage{scrlayer-scrpage}    % Manipulation des Seitenstils
%\pagestyle{scrheadings}  % Stil für die Seite setzen
%\automark{section}  % Abschnittsnamen als Seitenbeschriftung 
%\setheadsepline{.5pt}    % Kopzeile durch Linie abtrennen
\usepackage[hidelinks]{hyperref}  % Links und weitere PDF-Features
\usepackage{typearea, tikz, pgffor, verbatim}
\usetikzlibrary{calc, quotes}
\usepackage[cache=false, outputdir=build]{minted}
\usepackage{placeins} % provides \FloatBarrier

\definecolor{bg}{rgb}{0.9,0.9,0.9}
\newcommand\numberthis{\addtocounter{equation}{1}\tag{\theequation}}
\newcommand{\bra}{}
\newcommand{\kom}{}
\newcommand{\brak}{}
\newcommand{\braket}{}
\newcommand{\intx}{}
\newcommand{\intr}{}
\newcommand{\kla}{}
\newcommand{\klam}{}
\newcommand{\klammer}{}
\newcommand{\oper}{}
\newcommand{\tecxt}{}
\newcommand{\Bra}{}
\newcommand{\Ket}{}
\renewcommand{\i}{\text{i}}
\newcommand{\e}{\text{e}}
\newcommand*{\Fourier}{\textsc{Fourier}\ }
\newcommand*{\F}{\mathcal{F}}
\newcommand*{\sinc}{\text{sinc\,}}
\newcommand{\conv}{\mathop{\scalebox{3.5}{\raisebox{-0.38ex}{\makebox[0.5\width][c]{$\ast$}}}}\limits}

\newcommand*{\titel}{02 Spline Interpolation}
\newcommand*{\autor}{Joachim Schwardt und Julian Fleck}

\hypersetup{pdfauthor={\autor}, pdftitle={\titel}}

%\titlehead{titlehead}
\subject{Numerik I - Aufgaben}
\title{\titel}
%\author{\autor}
\author{\begin{tabular}{l|l}
Joachim Schwardt & 4768711 \\ 
Julian Fleck & 4759587\\ 
\end{tabular}}
%\author{\begin{tabular}{|l|l|}
%\hline
%Joachim Schwardt & 4768711 \\ \hline
%Julian Fleck & 4759587\\ \hline
%\end{tabular}}
\date{\today}

\begin{document}

%\begin{titlepage}
\maketitle  % Titel setzen
%\tableofcontents  % Inhaltsverzeichnis setzen
%\end{titlepage}

\section*{Aufgabe 1}
Sei $p_1(x,y)=ax^2+bxy+cy^2+dx+ey+f$ das Polynom, welches eine Funktion $f(x)$ an den $\mathbb{R}^3$-Punkten
\begin{align}
\kla\left( 0, 0, f(0,0) \right),&&&\kla\left( 0, \frac{1}{2},f\kla\left( 0,\frac{1}{2} \right) \right),&&\kla\left( 0,1,f(0,1) \right),\\
\kla\left( 1, 0, f(1,0) \right),&&&\kla\left( \frac{1}{2}, 0,f\kla\left( \frac{1}{2}, 0 \right) \right),&&\kla\left( \frac{1}{2},\frac{1}{2},f\kla\left( \frac{1}{2},\frac{1}{2} \right) \right)
\end{align}
interpoliert.
Das zu lösende Gleichungssystem lautet demnach
\begin{align}
f(0,0)&=p_1(0,0)=f,\\
f(1,0)&=p_1(1,0)=a+d+f,\\
f(\frac{1}{2},0)&=p_1(\frac{1}{2},0)=\frac{a}{4}+\frac{d}{2}+f,\\
f(0,1)&=p_1(0,1)=c+e+f,\\
f(0,\frac{1}{2})&=p_1(0,\frac{1}{2})=\frac{c}{4}+\frac{e}{2}+f,\\
f(\frac{1}{2},\frac{1}{2})&=p_1(\frac{1}{2},\frac{1}{2})=\frac{a+b+c}{4} + \frac{d+e}{2}+f.
\end{align}
Die Koeffizienten können iterativ bestimmt werden.
Man findet $f=f(0,0)$ und
\begin{align}
d &= 4f(\frac{1}{2},0)-f(1,0)-3f(0,0)&&\Rightarrow&a&=2f(0,0)+2f(1,0)-4f(\frac{1}{2},0).
\end{align}
Da das Vorgehen für $c,e$ analog zu dem für $a,d$ ist ($f(x,y)\rightarrow f(y,x)$,) kann man dann
\begin{align}
e &= 4f(0,\frac{1}{2})-f(0,1)-3f(0,0)&&\Rightarrow&c&=2f(0,0)+2f(0,1)-4f(0,\frac{1}{2}).
\end{align}
ablesen. 
Der letzte Koeffizient lautet dann
\begin{align}
b&=-a-c-2(d+e)-4f=4\kla\left( f(0,0)+f(\frac{1}{2},\frac{1}{2})-f(0,\frac{1}{2})-f(\frac{1}{2},0) \right).
\end{align}
Das vollständige Polynom ist also
\begin{align*}
p_1(x,y)&=2x^2\kla\left( f(0,0)+f(1,0)-2f(\frac{1}{2},0) \right)+2y^2\kla\left( f(0,0)+f(0,1)-2f(0,\frac{1}{2}) \right)+\\
&+x\kla\left( 4f(\frac{1}{2},0)-f(1,0)-3f(0,0) \right)+y\kla\left( 4f(0,\frac{1}{2})-f(0,1)-3f(0,0) \right)+\\
&+4xy\kla\left( f(0,0)+f(\frac{1}{2},\frac{1}{2})-f(0,\frac{1}{2})-f(\frac{1}{2},0) \right)+f(0,0). \numberthis\label{eqn:p1xy}
\end{align*}
Dasselbe Vorgehen führt natürlich auch zu $p_2$. 
Man kann sich aber überlegen, dass $p_2(x,y;f(x,y))\equiv p_1(1-x,1-y;f(1-x,1-y))$ ist.
Mit der Notation ist gemeint, dass auch die Stützstellen an den jeweilig verschobenen Koordinaten auszuwerten sind.
Damit erhält man dann direkt
\begin{align*}
p_2(x,y)&=2(1-x)^2\kla\left( f(1,1)+f(0,1)-2f(\frac{1}{2},1) \right)+2(1-y)^2\kla\left( f(1,1)+f(1,0)-2f(1,\frac{1}{2}) \right)+\\
&+(1-x)\kla\left( 4f(\frac{1}{2},1)-f(0,1)-3f(1,1) \right)+(1-y)\kla\left( 4f(1,\frac{1}{2})-f(1,0)-3f(1,1) \right)+\\
&+4(1-x)(1-y)\kla\left( f(1,1)+f(\frac{1}{2},\frac{1}{2})-f(1,\frac{1}{2})-f(\frac{1}{2},1) \right)+f(1,1). \numberthis\label{eqn:p2xy}
\end{align*}
Diese Form ist insbesondere hilfreich, wenn man nun im Folgenden die Polynome an den Punkten $(x,1-x)$ auswertet.
Für $p_1$ ergibt sich dann
\begin{align*}
p_1(x,1-x)&=f(0,0)\kla\left( 1+2x^2+2(1-x)^2-3x-3(1-x)+4x(1-x) \right) + \\
&+4x(1-x)f(\frac{1}{2},\frac{1}{2})+\\
&+f(1,0)\kla\left( 2x^2-x \right)+ \\
&+f(0,1)\kla\left( 2(1-x)^2-(1-x) \right)+\\
&+f(\frac{1}{2},0)\kla\left( -4x^2+4x-4x(1-x) \right)+\\
&+f(0,\frac{1}{2})\kla\left( -4(1-x)^2+4(1-x)-4(1-x)x \right)= \numberthis\\
&=4x(1-x)f(\frac{1}{2},\frac{1}{2})+f(1,0)\kla\left( 2x-1 \right)x+f(0,1)\kla\left( 2(1-x)-1 \right)(1-x).
\end{align*}
Aus Symmetriegründen gilt dann für $p_2$ wie zuvor
\begin{align}
p_2( x,&1-x;f(x,1-x) )=p_1\kla\left( 1-x, 1-(1-x); f(1-x, 1-(1-x)) \right)=\\
&=p_1(1-x,x;f(1-x,x))=\\
&=4(1-x)xf(\frac{1}{2},\frac{1}{2})+f(0,1)(2(1-x)-1)(1-x)+f(1,0)(2x-1)(1-x).
\end{align}
Damit ist die stückweise definierte Interpolationsfunktion also tatsächlich stetig.

\section*{Aufgabe 2}
\subsection*{a)}
Es gilt
\begin{align}
T_\alpha(s)&=\frac{\tecxt\text{d}\alpha(t(s))}{\tecxt\text{d}s}=\gamma'(t(s))t'(s)=\begin{pmatrix}
1 \\ g'(t(s))
\end{pmatrix}t'(s).
\end{align}
Aus der Bedingung $\|T_\alpha\|_2=1$ erhält man dann die gesuchte Differentialgleichung für $t(s)$ gemäß
\begin{align}
1&=\|T_\alpha\|_2^2=t'(s)^2\kla\left( 1+g'(t(s))^2 \right)&&\Rightarrow&t'(s)&=\frac{1}{\sqrt{1+g'(t(s))^2}}.
\end{align}
\subsection*{b)}
Aus der Bedingung
\begin{align}
0&=T_\alpha\cdot N_\alpha=t'(s)\kla\left( 1\cdot N_\alpha^{(x)} + g'(t(s))\cdot N_\alpha^{(y)} \right)
\end{align}
erhält man 
\begin{align}
N_\alpha^{(x)}&=-g'(t(s))N_\alpha^{(y)},
\end{align}
also
\begin{align}
N_\alpha&\propto\begin{pmatrix}
-g'(t(s)) \\ 1
\end{pmatrix}.
\end{align}
Über die Normierung kann man dann den Proportionalitätsfaktor $\beta$ bestimmen.
Man erhält also
\begin{align}
1&=\|N_\alpha\|_2^2=\beta^2\kla\left( g'(t(s))^2+1 \right)&&\Rightarrow&\beta&=\frac{1}{\sqrt{1+g'(t(s))^2}}=t'(s).
\end{align}
Insgesamt lautet der Normalenvektor
\begin{align}
N_\alpha(s)&=t'(s)\begin{pmatrix}
-g'(t(s)) \\ 1
\end{pmatrix}.
\end{align}
Für den nächsten Schritt ist es hilfreich, zunächst $T_\alpha'(s)$ zu berechnen, also
\begin{align}
T_\alpha'(s)&=\frac{\tecxt\text{d}}{\tecxt\text{d}s}\begin{pmatrix}
1 \\ g'(t(s))
\end{pmatrix}t'(s)=\begin{pmatrix}
1 \\ g'(t(s))
\end{pmatrix}t''(s)+\begin{pmatrix}
0 \\ g''(t(s))t'(s)
\end{pmatrix}t'(s)=\\
&=\begin{pmatrix}
t''(s) \\ t''(s)g'(t(s))+t'(s)^2g''(t(s))
\end{pmatrix}.
\end{align}
Es ist nicht offensichtlich, dass dieser Vektor parallel zu $N_\alpha$ ist.
Mit der Annahme $T'_\alpha(s)=k(s)N_\alpha(s)$ erhält man zunächst für die erste Komponente die Gleichung
\begin{align}
t''(s)&=-k(s)t'(s)g'(t(s))\\
\Rightarrow\ \ |k(s)|&=\Bigg\vert\frac{1}{-t'(s)g'(t(s))} \underbrace{\frac{-g'(t(s))g''(t(s))t'(s)}{\sqrt{1+g'(t(s))^2}^3}}_{=t''(s)} \Bigg\vert=\frac{|g''(t(s))|}{\kla\left( 1+g'(t(s))^2 \right)^{3/2}}.
\end{align}
Um die Proportionalität zu beweisen, muss auch die andere Komponente untersucht werden.
Man findet allerdings ebenfalls über
\begin{align}
k(s)t'(s)&=t''(s)g'(t(s))+t'(s)^2g''(t(s))\\
\Rightarrow\ \ |k(s)|&=\frac{1}{|t'(s)|}\left\vert  - \frac{g'(t(s))g''(t(s))t'(s)}{\sqrt{1+g'(t(s))^2}^3} g'(t(s))+t'(s)^2g''(t(s)) \right\vert=\\
&=|g''(t(s))t'(s)^3|\left\vert -g'(t(s))^2+\kla\left( 1+g'(t(s))^2 \right)\right\vert=|g''(t(s))t'(s)^3|,
\end{align}
was dem obigen Ergebnis entspricht.
\section*{Aufgabe 3}
Gegeben ist der Quadratur-Operator $Q[f]:=f(c)(b-a)$, mit welchem der Integralwert
\begin{align}
I&=\intx\int_{a}^{b}f(x)\, \mathrm{d}x
\end{align}
näherungsweise bestimmt werden soll.
\subsection*{a)}
Gemäß dem Mittelwertsatz existiert zu jedem $x\in[a,b]$ ein $\xi_x\in (a,b)$, sodass
\begin{align}
f(x)-f(c)&=f'(\xi_x)(x-c).
\end{align}
Allgemein gilt mit diesem Ansatz für den Fehler dann die Abschätzung
\begin{align}
|I-Q[f]|&=\left\vert\intx\int_{a}^{b}f(x)-f(c)\, \mathrm{d}x\right\vert\le\intx\int_{a}^{b}|f(x)-f(c)|\, \mathrm{d}x\le\\
&\le\max_{x\in[a,b]}\left\vert f'(x) \right\vert\intx\int_{a}^{b}|x-c|\, \mathrm{d}x.
\end{align}
Für $c=a$ folgt dann
\begin{align}
|I-Q[f]|&\le\intx\int_{0}^{b-a}x\, \mathrm{d}x \max_{x\in[a,b]}\left\vert f'(x) \right\vert=\frac{(b-a)^2}{2}\max_{x\in[a,b]}\left\vert f'(x) \right\vert.
\end{align}
Für $c=b$ findet man analog dasselbe Ergebnis, da
\begin{align}
\intx\int_{a}^{b}|x-a|\, \mathrm{d}x&=\intx\int_{0}^{b-a}x\, \mathrm{d}x\kom\overset{\substack{{u\,:=\,b-x} \\ |}}{=}\intx\int_{b}^{a}b-u\, -\mathrm{d}u=\intx\int_{a}^{b}|x-b|\, \mathrm{d}x.
\end{align}
\subsection*{b)}
Sei $c=\frac{a+b}{2}$ und $p_1(x)=f(c)+(x-c)f'(c)$.
Der Satz von Taylor besagt, dass der maximale Fehler dieses Taylorpolynoms vom Grad eins durch 
%http://pi.math.cornell.edu/~demlow/425/chap3.pdf
\begin{align}
\left\vert f(x)-p_1(x)\right\vert&\le \frac{1}{2!}\max_{x\in[a,b]}|f''(x)|(x-c)^2
\end{align}
beschrieben wird.
Ein kurzer Beweis dieser Ungleichung ist weiter unten gegeben.
Beachte, dass
\begin{align}
\intx\int_{a}^{b}x-c\, \mathrm{d}x&=\intx\int_{a}^{b}x-\frac{a+b}{2}\, \mathrm{d}x=\frac{b^2-a^2}{2}-\frac{a+b}{2}(b-a)=0. \label{eqn:int ab omega x}
\end{align}
Damit folgt die Fehlerabschätzung
\begin{align}
|I-Q[f]|&=\left\vert\intx\int_{a}^{b}f(x)-f(c)\, \mathrm{d}x \right\vert = \bigg\vert\intx\int_{a}^{b} f(x)-p_1(x)+\underbrace{\cancel{f'(c)(x-c)}}_{\eqref{eqn:int ab omega x}} \, \mathrm{d}x\bigg\vert\le\\
&\le\frac{1}{2}\max_{x\in[a,b]}\left\vert f''(x) \right\vert\intx\int_{a}^{b}\kla\left( x-\frac{a+b}{2} \right)^2\, \mathrm{d}x=\\
&=\frac{1}{2}\max_{x\in[a,b]}\left\vert f''(x) \right\vert 2\intx\int_{0}^{\frac{b-a}{2}}x^2\,\mathrm{d}x=\frac{(b-a)^3}{24}\max_{x\in[a,b]}\left\vert f''(x) \right\vert .
\end{align}
%Alternativ mit Sicherheit richtig, aber leider dafür falscher Vorfaktor:
%\begin{align}
%|I-Q[f]|&=\left\vert\intx\int_{a}^{b}f(x)-f(c)\, \mathrm{d}x \right\vert = \bigg\vert\intx\int_{a}^{b}\Big[ f'(\xi_x)-f'(c)+\underbrace{\cancel{f'(c)}}_{\eqref{eqn:int ab omega x}} \Big]\omega(x)\, \mathrm{d}x\bigg\vert\le\\
%&\le\max_{x\in[a,b]}\left\vert f'(x)-f'(c) \right\vert\intx\int_{a}^{b}\left\vert x-\frac{a+b}{2}\right\vert\, \mathrm{d}x\le\\
%&\le\max_{x\in[a,b]}\left\vert f''(x) \right\vert \max_{x\in[a,b]}\left\vert x-c\right\vert 2\intx\int_{0}^{\frac{b-a}{2}}x\, \mathrm{d}x=\\
%&=\frac{b-a}{2}\max_{x\in[a,b]}\left\vert f''(x)\right\vert \kla\left( \frac{b-a}{2} \right)^2\le \frac{(b-a)^3}{8}\max_{x\in[a,b]}\left\vert f''(x)\right\vert.
%\end{align}
%\begin{align}
%|I-Q[f]|&=\left\vert\intx\int_{a}^{b}f(x)-f(c)\, \mathrm{d}x \right\vert = \bigg\vert\intx\int_{a}^{b}\Big[ f'(\xi_x)-f'(c)+\underbrace{\cancel{f'(c)}}_{\eqref{eqn:int ab omega x}} \Big]\omega(x)\, \mathrm{d}x\bigg\vert\le\\
%&\le\max_{x\in[a,b]}\left\vert \frac{f'(\xi_x)-f'(c)}{x-c} \right\vert\intx\int_{a}^{b}\left( x-c\right)^2\, \mathrm{d}x\le\\
%&\le\max_{x\in[a,b]}\left\vert f''(x) \right\vert 2\intx\int_{0}^{\frac{b-a}{2}}x^2\, \mathrm{d}x=\\
%&=\frac{b-a}{2}\max_{x\in[a,b]}\left\vert f''(x)\right\vert \kla\left( \frac{b-a}{2} \right)^2\le \frac{(b-a)^3}{24}\max_{x\in[a,b]}\left\vert f''(x)\right\vert.
%\end{align}


%%% !!! IMPORTANT RESOURCE !!! %%%
%%% https://math.stackexchange.com/questions/481661/simplest-proof-of-taylors-theorem
Hier noch ein kurzer Beweis für den verwendeten Fall des Satzes von Taylor.
Schreibe zunächst eine beliebige Funktion $f\in\mathcal{C}^2$ mittels des Hauptsatzes der Differential- und Integralrechnung iterativ um.
Also 
\begin{align}
f(x)&=f(x_0)+\intx\int_{x_0}^{x}f'(t_1)\, \mathrm{d}t_1=f(x_0)+\intx\int_{x_0}^{x}f'(x_0)+\intx\int_{x_0}^{t_1}f''(t_2)\, \mathrm{d}t_2\, \mathrm{d}t_1=\\
&=f(x_0)+f'(x_0)(x-x_0)+\intx\int_{x_0}^{x}\intx\int_{x_0}^{t_1}f''(t_2)\, \mathrm{d}t_2\, \mathrm{d}t_1
\end{align}
Dann gilt für den Fehler
\begin{align}
\left\vert f(x)-f(x_0)-f'(x_0)(x-x_0)\right\vert&\le \max_{t\in [x_0,x]}\left\vert f''(t)\right\vert \intx\int_{x_0}^{x}\intx\int_{x_0}^{t_1}\, \mathrm{d}t_2\, \mathrm{d}t_1=\\
&=\frac{1}{2}(x-x_0)^2\max_{t\in [x_0,x]}\left\vert f''(t)\right\vert.
\end{align}
In unserem Fall ist $x_0=\frac{a+b}{2}$ und $x\in [a,b]$.
Das Maximum kann natürlich auch über $[a,b]$ ausgewertet werden, da das Intervall $[x_0,x]$ für alle $x$ darin enthalten ist.

%Für $c=\frac{a+b}{2}$ kann man das Integral aus der Fehlerabschätzung aufteilen.
%Schreibe
%\begin{align}
%\intx\int_{a}^{b}f(x)-f(c)\, \mathrm{d}x&=\intx\int_{a}^{\frac{a+b}{2}}f(x)-f(c)\, \mathrm{d}x+\intx\int_{\frac{a+b}{2}}^{b}f(x)-f(c)\, \mathrm{d}x=\\
%&=\intx\int_{a}^{\frac{a+b}{2}}-|f(x)-f(c)|\, \mathrm{d}x+\intx\int_{\frac{a+b}{2}}^{b}|f(x)-f(c)|\, \mathrm{d}x.
%\end{align}
%Für beide Teilintegrale kann nun die Fehlerabschätzung aus a) verwendet werden.
%Daraus folgt
%\begin{align}
%|I-Q[f]|&\le \left\vert -\frac{\kla\left( \frac{a+b}{2}-a \right)^2}{2}\max_{x\in[a,\frac{a+b}{2}]}\left\vert f'(x) \right\vert + \frac{\kla\left( b-\frac{a+b}{2} \right)^2}{2}\max_{x\in[\frac{a+b}{2},b]}\left\vert f'(x) \right\vert \right\vert\le\\
%&\le\frac{(b-a)^2}{8}\left\vert\max_{x\in[\frac{a+b}{2},b]}\left\vert f'(x) - f'(c) \right\vert - \max_{x\in[a,\frac{a+b}{2}]}\left\vert f'(x) - f'(c) \right\vert\right\vert.
%\end{align}
%Es gilt
%\begin{align}
%f'(x)&=f''(\xi_x)(x-c)+f'(c)\\
%\Rightarrow\ \ \max_{x\in[a,\frac{a+b}{2}]}|f'(x)|&\le |f'(c)|+\max_{x\in[a,\frac{a+b}{2}]}|f''(x)|
%\end{align}
%Für die Differenzen der ersten Ableitungen gilt wieder der Mittelwertsatz, also
%\begin{align}
%\max_{x\in[a,\frac{a+b}{2}]}\left\vert f'(x) - f'(c) \right\vert&\le\max_{x\in[a,\frac{a+b}{2}]}\left\vert f''(x)\left\vert x-\frac{a+b}{2} \right\vert\right\vert =\frac{b-a}{2}\max_{x\in[a,\frac{a+b}{2}]}\left\vert f''(x)\right\vert
%\end{align}
%und analog
%\begin{align}
%\max_{x\in[\frac{a+b}{2},b]}\left\vert f'(x) - f'(c) \right\vert&\le\max_{x\in[\frac{a+b}{2},b]}\left\vert f''(x)\left\vert x-\frac{a+b}{2} \right\vert\right\vert =\frac{b-a}{2}\max_{x\in[\frac{a+b}{2},b]}\left\vert f''(x)\right\vert.
%\end{align}
%Man benötigt noch die Ungleichung $|x-y|\le |x|+|y|$
%\begin{thebibliography}{9}
%\bibitem{example} description, \url{https://www.exmapleurl}, day.\,month~2021.
%\end{thebibliography}

\end{document}